% Part: counterfactuals
% Chapter: minimal-change-semantics
% Section: introduction

\documentclass[../../../include/open-logic-section]{subfiles}

\begin{document}

\olfileid[pt]{cnt}{min}{int}

\olsection{Introdução}

Stalnaker e Lewis propuseram explicações de condicionais contrafactuais
como ``Se o fósforo fosse riscado, ele acenderia''. Suas explicações
eram propostas de como compreender corretamente as condições de verdade de
tais sentenças. A ideia subjacente a ambas é esta: para avaliar
se um condicional contrafactual é verdadeiro, precisamos considerar
os mundos possíveis que diferem minimamente do modo como o
mundo real é, de maneira a tornar verdadeiro o antecedente. Se o consequente é
verdadeiro nesses mundos possíveis, então o contrafactual é verdadeiro. Por
exemplo, suponha que eu tenha um fósforo e uma caixa de fósforos na mão. No
mundo real, apenas olho para eles e pondero o que aconteceria se eu
riscasse o fósforo. A mudança mínima em relação ao mundo real
na qual risco o fósforo é aquela em que decido agir e riscá-lo.
Ela é mínima porque nada mais muda: não pulo também
no ar, riscar o fósforo não incendeia também meu cabelo, não
perco de repente toda a força dos dedos, não sou simultaneamente
encharcado em uma emboscada com uma Super Soaker etc. Nessa possibilidade
alternativa, o fósforo acende. Portanto, é verdade que, se eu
riscasse o fósforo, ele acenderia.

Essa explicação intuitiva pode ser combinada com semânticas formais para lógicas
de contrafactuais. Lewis introduziu o símbolo ``$\boxright$'' para
o contrafactual, enquanto Stalnaker usou o símbolo~``$>$''. Usaremos
~$\cif$ e o acrescentaremos como conectivo binário à lógica
proposicional. Assim, além de !!{formula}s da forma $!A \lif
!B$, também temos !!{formula}s da forma~$!A \cif !B$. A semântica formal,
como a semântica relacional da lógica modal, baseia-se em modelos nos
quais as !!{formula}s são avaliadas em mundos, e a condição de satisfação
que define $\mSat{M}{!A \cif !B}[w]$ é dada em termos de
$\mSat{M}{!A}[w']$ e $\mSat{M}{!B}[w']$ para alguns outros
mundos~$w'$. Quais $w'$? Intuitivamente, aquele ou aqueles mais próximos de~$w$
para os quais vale~$\mSat{M}{!A}[w']$. Isso exige que uma relação
de ``proximidade'' também seja incluída no modelo.

Lewis introduziu uma maneira instrutiva de representar graficamente situações
contrafactuais. Cada mundo possível está no centro de um conjunto
de esferas aninhadas que contêm outros mundos---desenhamos essas esferas como
círculos concêntricos. Os mundos entre duas esferas são igualmente próximos
do mundo central; os contidos em uma esfera mais interna são mais próximos,
e os que estão em uma esfera circundante são mais distantes.
\begin{center}
\begin{tikzpicture}[scale=.7]\small
  \spheresystem{5}
  \draw (0,0) node {$w$};
  \propositionintersect{0}{4}{40}{3.5}
  {\tikzset{proposition/.append={smooth,tension=1.4}}
  \proposition[shift={(1.6,-1)}]{90}{1}{30}{4}}
  \path (1.5,3) node[below] {$\formula{B}$};
  \path (3,0) node {$\formula{A}$};
\end{tikzpicture}
\end{center}
Os mundos-$!A$ mais próximos são os mundos $w'$ nos quais~$!A$ é satisfeita
e que estão na menor esfera ao redor do mundo central~$w$ (a área
cinza). Intuitivamente, $!A \cif !B$ é satisfeita em~$w$ se $!B$ é verdadeira
em todos os mundos-$!A$ mais próximos.

\end{document}
